% Options for packages loaded elsewhere
\PassOptionsToPackage{unicode}{hyperref}
\PassOptionsToPackage{hyphens}{url}
%
\documentclass[
  letterpaper,
]{book}

\usepackage{amsmath,amssymb}
\usepackage{iftex}
\ifPDFTeX
  \usepackage[T1]{fontenc}
  \usepackage[utf8]{inputenc}
  \usepackage{textcomp} % provide euro and other symbols
\else % if luatex or xetex
  \usepackage{unicode-math}
  \defaultfontfeatures{Scale=MatchLowercase}
  \defaultfontfeatures[\rmfamily]{Ligatures=TeX,Scale=1}
\fi
\usepackage{lmodern}
\ifPDFTeX\else  
    % xetex/luatex font selection
\fi
% Use upquote if available, for straight quotes in verbatim environments
\IfFileExists{upquote.sty}{\usepackage{upquote}}{}
\IfFileExists{microtype.sty}{% use microtype if available
  \usepackage[]{microtype}
  \UseMicrotypeSet[protrusion]{basicmath} % disable protrusion for tt fonts
}{}
\makeatletter
\@ifundefined{KOMAClassName}{% if non-KOMA class
  \IfFileExists{parskip.sty}{%
    \usepackage{parskip}
  }{% else
    \setlength{\parindent}{0pt}
    \setlength{\parskip}{6pt plus 2pt minus 1pt}}
}{% if KOMA class
  \KOMAoptions{parskip=half}}
\makeatother
\usepackage{xcolor}
\setlength{\emergencystretch}{3em} % prevent overfull lines
\setcounter{secnumdepth}{5}
% Make \paragraph and \subparagraph free-standing
\makeatletter
\ifx\paragraph\undefined\else
  \let\oldparagraph\paragraph
  \renewcommand{\paragraph}{
    \@ifstar
      \xxxParagraphStar
      \xxxParagraphNoStar
  }
  \newcommand{\xxxParagraphStar}[1]{\oldparagraph*{#1}\mbox{}}
  \newcommand{\xxxParagraphNoStar}[1]{\oldparagraph{#1}\mbox{}}
\fi
\ifx\subparagraph\undefined\else
  \let\oldsubparagraph\subparagraph
  \renewcommand{\subparagraph}{
    \@ifstar
      \xxxSubParagraphStar
      \xxxSubParagraphNoStar
  }
  \newcommand{\xxxSubParagraphStar}[1]{\oldsubparagraph*{#1}\mbox{}}
  \newcommand{\xxxSubParagraphNoStar}[1]{\oldsubparagraph{#1}\mbox{}}
\fi
\makeatother


\providecommand{\tightlist}{%
  \setlength{\itemsep}{0pt}\setlength{\parskip}{0pt}}\usepackage{longtable,booktabs,array}
\usepackage{calc} % for calculating minipage widths
% Correct order of tables after \paragraph or \subparagraph
\usepackage{etoolbox}
\makeatletter
\patchcmd\longtable{\par}{\if@noskipsec\mbox{}\fi\par}{}{}
\makeatother
% Allow footnotes in longtable head/foot
\IfFileExists{footnotehyper.sty}{\usepackage{footnotehyper}}{\usepackage{footnote}}
\makesavenoteenv{longtable}
\usepackage{graphicx}
\makeatletter
\newsavebox\pandoc@box
\newcommand*\pandocbounded[1]{% scales image to fit in text height/width
  \sbox\pandoc@box{#1}%
  \Gscale@div\@tempa{\textheight}{\dimexpr\ht\pandoc@box+\dp\pandoc@box\relax}%
  \Gscale@div\@tempb{\linewidth}{\wd\pandoc@box}%
  \ifdim\@tempb\p@<\@tempa\p@\let\@tempa\@tempb\fi% select the smaller of both
  \ifdim\@tempa\p@<\p@\scalebox{\@tempa}{\usebox\pandoc@box}%
  \else\usebox{\pandoc@box}%
  \fi%
}
% Set default figure placement to htbp
\def\fps@figure{htbp}
\makeatother

\usepackage{booktabs}
\usepackage{lscape}

\makeatletter
\@ifpackageloaded{bookmark}{}{\usepackage{bookmark}}
\makeatother
\makeatletter
\@ifpackageloaded{caption}{}{\usepackage{caption}}
\AtBeginDocument{%
\ifdefined\contentsname
  \renewcommand*\contentsname{Daftar Isi}
\else
  \newcommand\contentsname{Daftar Isi}
\fi
\ifdefined\listfigurename
  \renewcommand*\listfigurename{Daftar Gambar}
\else
  \newcommand\listfigurename{Daftar Gambar}
\fi
\ifdefined\listtablename
  \renewcommand*\listtablename{Daftar Tabel}
\else
  \newcommand\listtablename{Daftar Tabel}
\fi
\ifdefined\figurename
  \renewcommand*\figurename{Gambar}
\else
  \newcommand\figurename{Gambar}
\fi
\ifdefined\tablename
  \renewcommand*\tablename{Tabel}
\else
  \newcommand\tablename{Tabel}
\fi
}
\@ifpackageloaded{float}{}{\usepackage{float}}
\floatstyle{ruled}
\@ifundefined{c@chapter}{\newfloat{codelisting}{h}{lop}}{\newfloat{codelisting}{h}{lop}[chapter]}
\floatname{codelisting}{Daftar}
\newcommand*\listoflistings{\listof{codelisting}{Daftar Daftar}}
\makeatother
\makeatletter
\makeatother
\makeatletter
\@ifpackageloaded{caption}{}{\usepackage{caption}}
\@ifpackageloaded{subcaption}{}{\usepackage{subcaption}}
\makeatother

\ifLuaTeX
\usepackage[bidi=basic]{babel}
\else
\usepackage[bidi=default]{babel}
\fi
\babelprovide[main,import]{indonesian}
% get rid of language-specific shorthands (see #6817):
\let\LanguageShortHands\languageshorthands
\def\languageshorthands#1{}
\usepackage{bookmark}

\IfFileExists{xurl.sty}{\usepackage{xurl}}{} % add URL line breaks if available
\urlstyle{same} % disable monospaced font for URLs
\hypersetup{
  pdftitle={Quarto Template},
  pdfauthor={Bakti Siregar, M.Sc.},
  pdflang={id},
  hidelinks,
  pdfcreator={LaTeX via pandoc}}


\title{Quarto Template}
\author{Bakti Siregar, M.Sc.}
\date{2026-11-22}

\begin{document}
\frontmatter
\maketitle

\renewcommand*\contentsname{Daftar Isi}
{
\setcounter{tocdepth}{2}
\tableofcontents
}

\mainmatter
\bookmarksetup{startatroot}

\chapter{Pengenalan}\label{pengenalan}

\bookmarksetup{startatroot}

\chapter*{Kata Pengantar}\label{kata-pengantar}
\addcontentsline{toc}{chapter}{Kata Pengantar}

\markboth{Kata Pengantar}{Kata Pengantar}

Dengan mengucap syukur kepada Tuhan Yang Maha Esa, saya sangat senang
dapat mempersembahkan eBook ``Quarto Template'' kepada para
pembaca/pengguna.

\section*{Penulis}\label{penulis}
\addcontentsline{toc}{section}{Penulis}

\markright{Penulis}

\begin{itemize}
\tightlist
\item
  \textbf{\href{https://www.linkedin.com/in/dsciencelabs/}{Bakti
  Siregar, M.Sc., CDS}} bekerja sebagai Dosen di Prodi Sains Data
  Institut Teknologi Sains Bandung. Beliau adalah meraih gelar
  Magister-nya dari Departemen Matematika Terapan (Applied Mathematics)
  National Sun Yat Sen University, Taiwan. Selain mengajar beliau juga
  pernah bekerja sebagai \textbf{Data Scientist Freelance} di
  perusahaan-perusahaan ternama seperti
  \href{https://www.jne.co.id/id/beranda}{JNE},
  \href{https://www.samoragroup.co.id/home/en}{Samora Group},
  \href{https://www.pertamina.com/}{Pertamina}, dan saat ini masih aktif
  sebagai \textbf{Data Scientist Freelance} di
  \href{PT.\%20Green\%20City\%20Traffic}{PT. Green City Traffic}. Beliau
  memiliki antusiasme khusus dalam mengerjakan proyek (mengajar) Big
  Data Analytics, Machine Learning, Optimisasi, dan Analisis Time Series
  di bidang keuangan dan investasi. Keahlian utama yang dimilikinya
  adalah bahasa pemrograman Statistik seperti R Studio dan Python.
  Beliau juga sudah terbiasa dalam mengaplikasikan sistem basis data
  MySQL/NoSQL untuk manajemen data, serta mahir dalam menggunakan tools
  Big Data seperti Spark dan Hadoop. Beberapa project beliau dapat
  dilihat di link berikut: \href{https://rpubs.com/dsciencelabs}{Rpubs},
  \href{https://github.com/dsciencelabs}{Github},
  \href{https://dsciencelabs.github.io/web/index.html}{Website}, dan
  \href{https://www.kaggle.com/baktisiregar/code}{Kaggle}.
\end{itemize}

\section*{Ucapan Terima Kasih}\label{ucapan-terima-kasih}
\addcontentsline{toc}{section}{Ucapan Terima Kasih}

\markright{Ucapan Terima Kasih}

Proses penulisan Tempate eBook ini tidak terlepas dari dukungan berbagai
pihak. Saya ingin mengucapkan terima kasih kepada:

\begin{itemize}
\tightlist
\item
  \textbf{Keluarga} yang selalu memberikan dukungan moral dan semangat
  tanpa henti.
\item
  \textbf{Rekan-rekan dan Kolaborator} yang telah memberikan masukan,
  saran, dan kritik yang konstruktif.
\item
  Institusi dan Organisasi, Khususnya
  \textbf{\href{https://itsb.ac.id/}{ITSB}} yang telah menyediakan
  sumber daya dan fasilitas yang diperlukan selama proses penelitian dan
  penulisan.
\end{itemize}

Saya berharap Template eBook ini dapat menjadi referensi yang bermanfaat
dan memberikan inspirasi serta pengetahuan baru bagi pembaca. Semoga
Ebook ini dapat memenuhi ekspektasi dan kebutuhan para pembaca dan
semoga ilmu yang disampaikan dapat bermanfaat bagi semua.

\section*{Masukan \& Saran}\label{masukan-saran}
\addcontentsline{toc}{section}{Masukan \& Saran}

\markright{Masukan \& Saran}

Semua masukan dan tanggapan Anda sangat berarti bagi kami untuk
memperbaiki Ebook Quarto Template ini kedepannya. Bagi para
pembaca/pengguna yang ingin menyampaikan masukan dan tanggapan,
dipersilahkan melalui kontak dibawah ini!

\textbf{Email:} dsciencelabs@outlook.com

\bookmarksetup{startatroot}

\chapter{Bab1}\label{bab1}

\section{Praktikum}\label{praktikum}

\bookmarksetup{startatroot}

\chapter{Bab2}\label{bab2}

\section{Praktikum}\label{praktikum-1}

\bookmarksetup{startatroot}

\chapter{Bab3}\label{bab3}

\section{Praktikum}\label{praktikum-2}

\bookmarksetup{startatroot}

\chapter{Referensi}\label{referensi}


\backmatter


\end{document}
